\chapter{Gates, controls and tests}

\label{ch:gates}

This chapter answers what the gate actually runs, what a green result does and does not prove,
and why so much of the surrounding machinery is about the \emph{shape} of a check rather than
its subject. The gate is one list of twenty steps, executed in two different trees, with the
environment sanitised so that a result cannot depend on who started it. Around it sit the
rules that decide whether a check is worth having: a criterion that passes before the work
measures nothing, a rule that suppresses needs a control that must still take the normal path,
and a check must consume what the live path produces rather than a copy of it.

\section{One list, read twice}

\texttt{gate-steps.txt} is the gate. Its header says so, and the format is one line per step:

\begin{lstlisting}[language=bash]
# Format: <label>|<shell command>   (one line per step; # lines and blank lines are ignored)
# A step passes only if its command exits 0; every step runs even after an earlier failure.
fmt|cargo fmt --all --check
clippy|CARGO_TARGET_DIR=".../clippy-gate" cargo clippy --workspace --all-targets -- -D warnings
build|CARGO_TARGET_DIR=".../build-gate" cargo build --workspace && cp ... target/debug/pii
lint-files|CARGO_TARGET_DIR="_tmp/tools-gate" cargo run -q -p xtask -- lint-files
\end{lstlisting}

\texttt{quality-gate.sh} reads that file for both modes; nothing else lists the steps. The
local run iterates the file, calls \texttt{run\_step} for each line, and keeps a status that
any failure sets.

\begin{lstlisting}[language=bash]
while IFS= read -r line || [ -n "$line" ]; do
    case "$line" in '' | '#'*) continue ;; esac
    run_step "${line%%|*}" "${line#*|}" || status=1
    steps=$((steps + 1))
done <"$GATE_STEPS"
# An empty or unreadable list must never read as green.
if [ "$steps" -eq 0 ]; then
    echo "FAIL: gate-steps.txt lists no steps"
    status=1
fi
\end{lstlisting}

Two properties are in that excerpt. Every step runs even after an earlier failure, so one run
produces the whole list of problems rather than the first one. And an empty list is a failure,
not a pass: a gate that runs nothing and exits zero is the most dangerous possible
configuration, because it looks exactly like a gate that ran everything.

\section{The twenty steps}

The list has grown by addition, and each addition has a reason recorded next to it in the
file. The table below is the inventory: what each step runs, and the class of defect it exists
to catch.

\begin{longtable}{@{}p{0.24\textwidth}p{0.34\textwidth}p{0.34\textwidth}@{}}
\toprule
\textbf{Step} & \textbf{Runs} & \textbf{Why it exists} \\
\midrule
\endhead
\texttt{fmt} & \texttt{cargo fmt --all --check} & zero-tolerance formatting \\
\texttt{clippy} & \texttt{clippy --workspace --all-targets -D warnings}, own target dir & one
\texttt{--all-targets} rlib hash can poison the doctest linker (E0463) \\
\texttt{clippy (hal-gpiod)} & \texttt{clippy -p pii-hal --features hal-gpiod}, own target dir & the
real GPIO backend is feature-gated off everywhere else; compile and lint only, never run \\
\texttt{build} & \texttt{cargo build --workspace}, then publish \texttt{target/debug/pii} & the
worker's supervisor tests spawn the real binary at a fixed path \\
\texttt{test (lib/bins/tests)} & \texttt{cargo test --workspace --lib --bins --tests} & the suite \\
\texttt{test (doc)} & \texttt{cargo test --workspace --doc}, own target dir & doctests build libs
without \texttt{cfg(test)}; sharing a target dir picks a stale rlib \\
\texttt{coverage-floor} & \texttt{cargo llvm-cov} then \texttt{xtask coverage-floor} & line coverage
must stay at or above \texttt{coverage-floor.txt} \\
\texttt{lint-files} & \texttt{xtask lint-files} & 450 effective lines per file, plus spec/doc
hygiene \\
\texttt{lint-scratch} & \texttt{xtask lint-scratch} & a test that uses a crate-local
\texttt{\_tmp} it
never creates \\
\texttt{lint-observability} & \texttt{xtask lint-observability} & every crate uses the
\texttt{pii-observe} seam, no new raw prints, registry names in \texttt{ALL} \\
\texttt{reachability-check} & \texttt{xtask reachability-check} & a declared capability seam has
a real
non-test caller \\
\texttt{deps-check} & \texttt{xtask deps-check} & five dependency invariants, from \texttt{cargo
metadata} \\
\texttt{net-check} & \texttt{xtask net-check} & sockets and name resolution only in allowlisted
files \\
\texttt{doc-drift-check} & \texttt{xtask doc-drift-check} & every backticked repo path in
\texttt{HANDOFF.md} exists \\
\texttt{autonomy (shell tests)} & \texttt{bash scripts/autonomy/run-tests.sh} & tests \texttt{cargo
test} cannot reach; without this the gate cannot go red on them \\
\texttt{plan-verify} & \texttt{sync-verify.sh --check} & an open row's verify column must still
be its
spec's criteria \\
\texttt{item-ids} & \texttt{item-ids.sh --check} & a flow item's id is its work's spec number \\
\texttt{devbox-check} & \texttt{modules/devbox/check.sh} & the devbox module's offline claims
hold \\
\texttt{install-config} & \texttt{deploy/install-config-keeps.test.sh} & a re-install keeps the
operator's config \\
\texttt{daemon-guard} & \texttt{scripts/checks/daemon-tests-are-guarded.sh} & a test that spawns a
live daemon owns it through \texttt{DaemonGuard} \\
\bottomrule
\end{longtable}

The autonomy step's history is the clearest argument for the list being a file rather than a
habit. \texttt{tick-message-worker-rc.test.sh} was red for two days after a change to
\texttt{tick.sh}'s phone line, and every gate in that window was green, because \texttt{cargo
test} never reached a shell test. A check that is not in the list does not exist.

\begin{center}
\resizebox{\linewidth}{!}{%
\begin{tikzpicture}[node distance=8mm and 12mm]
	\node[b] (s1) {fmt\\\scriptsize step 1};
	\node[b, right=of s1] (s2) {clippy\\\scriptsize steps 2--3};
	\node[b, right=of s2] (s3) {build\\\scriptsize step 4};
	\node[b, right=of s3] (s4) {tests\\\scriptsize steps 5--6};
	\node[b, below=of s4] (s5) {measure and lint\\\scriptsize steps 7--14};
	\node[b, left=of s5] (s6) {queue and host guards\\\scriptsize steps 15--20};

	\draw[f] (s1) -- (s2);
	\draw[f] (s2) -- (s3);
	\draw[f] (s3) -- (s4);
	\draw[f] (s4) -- (s5);
	\draw[f] (s5) -- (s6);
\end{tikzpicture}}
\captionof{figure}{The twenty steps as five stages and a tail of guards, where the numbered
groups are the ones a reader will otherwise have to count in the file.}
\end{center}

\section{Two trees}

The same list runs in two places, and the difference is what the result proves.

\begin{center}
\resizebox{\linewidth}{!}{%
\begin{tikzpicture}[node distance=8mm and 12mm]
	\node[b] (repo) {working tree};
	\node[b, below=of repo] (archive) {\texttt{git archive HEAD}};
	\node[b, right=of archive] (export) {export dir\\\scriptsize \texttt{fresh-gate-<pid>}};
	\node[b, above=of export] (overlay) {overlay\\\scriptsize modified $\cdot$ untracked $\cdot$ deletions};
	\node[b, right=of export] (gate) {\texttt{quality-gate.sh}\\\scriptsize 20 steps};
	\node[v, below=of gate] (pass) {pass?};
	\node[b, left=of pass] (kept) {kept on failure\\\scriptsize \texttt{gate-output.log}};
	\node[b, right=of pass] (gone) {removed on pass};

	\draw[f] (repo) -- (archive);
	\draw[f] (archive) -- (export);
	\draw[f] (overlay) -- (export);
	\draw[f] (export) -- (gate);
	\draw[f] (gate) -- (pass);
	\draw[f] (pass) -- node[p, above=2pt, midway]{no} (kept);
	\draw[f] (pass) -- node[p, above=2pt, midway]{yes} (gone);
\end{tikzpicture}}
\captionof{figure}{When the two trees disagree the export is the one to believe, because tests
that pass only because an earlier run left a directory behind are the class the second tree
catches.}
\end{center}

\texttt{quality-gate.sh} has three entry points: \texttt{--routine} (default, this tree, warm
cache, minutes), \texttt{--fresh} (the export, the release boundary), and \texttt{--summary}
(prints what the gate cannot exercise and runs nothing). The fresh mode is a wrapper around
\texttt{cargo xtask fresh-gate --worktree}, with a separate target dir, because an export's
artifacts must not land in the warm cache: a test binary bakes its manifest directory at
compile time, and a routine run that reused an artifact from a deleted export failed three
tests on a tree where the package's own suite passed 5/5.

\paragraph{Overlay, including deletion.} \texttt{overlay\_worktree} copies every modified and
untracked path from the working tree into the export, and then \emph{removes} from the export
every path git reports as deleted. The second half was added after a change that removes a
test file could not be gated at all: the export still held the deleted file while the overlaid
manifest no longer declared its dependency, so the gate failed on a tree that does not exist.
A gate that cannot express a deletion is not testing the change.

\paragraph{Single-file substitution.} \texttt{--overlay src=dest} copies one file over the
export. \texttt{dest} must be relative and must not contain \texttt{..}; both are refused at
argument parsing, before anything is created. This is the way to test a deliberately mutated
file without touching the tree.

\paragraph{Which list ran.} \texttt{check\_step\_list} reads \texttt{GATE\_STEPS}, counts the
non-comment lines in the \emph{export}, and refuses a reduced list unless \texttt{--quick} was
asked for:

\begin{lstlisting}[language=bash]
GATE_STEPS=gate-steps-quick.txt cargo xtask fresh-gate --worktree
# REFUSED: GATE_STEPS=gate-steps-quick.txt is the ITEM list (5 steps), not the gate that decides
# shipping. Unset GATE_STEPS (or set it to gate-steps.txt) for the full list, or pass --quick if that
# reduced run IS what you intend - it will say so in its own output.
\end{lstlisting}

That refusal exists because of a measured evening: an autonomy script exports the quick list
for the item boundary, a hand run that sourced it reported ``gated'' five times, and the
commit gate --- which runs the full list --- had three real failures waiting. The quick list
is a legitimate tool. It must not be the accidental default, so it is allowed only when asked
for, and the run always prints which list it used and how many steps it holds.

\section{The environment must not be inherited}

A gate that depends on the caller's environment is not a gate. Two variables are removed from
the child and one is redirected:

\begin{itemize}
	\item \texttt{PII\_STATE\_DIR} and \texttt{PII\_RUNTIME\_DIR} are removed. Measured: a lane
	daemon holding the real store made a daemon test fail with \texttt{channel.redb is locked:
	another daemon holds it}, because the test inherited the state dir. The test's own store
	parent was not enough --- the durable stores follow the variable.
	\item \texttt{CARGO\_TARGET\_DIR} points at the warm cache, so the export's tree is fresh
	and its compiler output is not. Pinning it inside the export made every attempt pay a full
	cold workspace build: a spec that took minutes took 40--60 minutes after the export entered
	the loop.
	\item \texttt{TEST\_GATE\_TARGET} points \emph{inside} the export for the test step only,
	because a test binary bakes \texttt{CARGO\_MANIFEST\_DIR} and a binary left in the shared
	cache would be reused by the next run.
\end{itemize}

The third is the subtle one, and it has its own test with a control: the test target must
start with the export's path, must not equal the shared cache's test directory, and every
other step must still get the warm cache. Without the control, a change that redirected
\emph{every} step would pass.

\section{Disk, retention and the two ceilings}

A fresh export is about 8 GB before it builds anything, and a coverage-measuring export needs
at least 11.4 GiB free. Both numbers are measurements, not estimates: a run with 11 GB free
drove the volume to zero for about three minutes and left an 11.4 GB failed export behind. The
floor is \texttt{FRESH\_GATE\_MIN\_FREE\_MB}, default 12\,288, checked before the export
directory is created, and the refusal names the free megabytes and the override.

The warm cache has the opposite problem: a bigger cache builds faster, so nothing bounded it
and it reached 105.8 GiB. It is now capped by \texttt{FRESH\_GATE\_CACHE\_MAX\_GB} (default
40) and trimmed before the build --- but never while a cargo or rustc is live, because a
running build would lose its target underneath it. When a build is live the tool says so and
skips the trim rather than guessing.

\paragraph{Retention by liveness.} Failed exports are removed before a new one is created, but
the predicate is not age. Each export directory is named for the pid that created it
(\texttt{fresh-gate-<pid>}), and retention asks whether that pid still exists. The measurement
behind this is worth repeating: retention once removed a \emph{running} gate's export, and
that gate then failed sixteen steps on \texttt{target/gate-steps/<step>.log: No such file or
directory} with \texttt{failing tests: []} --- a red verdict blaming the tree for a directory
another process deleted.

Two tests pin the pair: a running owner's export is not reaped, and --- as the control --- a
dead owner's export is. The control matters because ``never remove anything'' would also pass
the first test.

\begin{center}
\resizebox{\linewidth}{!}{%
\begin{tikzpicture}[node distance=8mm and 12mm]
	\node[b] (dir) {\texttt{fresh-gate-<pid>}};
	\node[v, right=of dir] (safe) {name ok?};
	\node[v, below=of safe] (alive) {alive?};
	\node[b, right=of alive, yshift=6mm] (keep) {keep};
	\node[b, right=of alive, yshift=-6mm] (reap) {remove};
	\node[b, left=of alive, xshift=-4mm] (skip) {leave alone\\\scriptsize not ours};

	\draw[f] (dir) -- (safe);
	\draw[f] (safe) -- node[p, right=2pt, midway]{yes} (alive);
	\draw[f] (alive) -- node[p, above=2pt, midway]{yes} (keep);
	\draw[f] (alive) -- node[p, below=2pt, midway]{no} (reap);
	\draw[bk] (safe.west) -- node[p, above=2pt, midway]{no} (skip);
\end{tikzpicture}}
\captionof{figure}{The name guard runs first and fails safe: an empty base, a foreign
directory or a bare prefix is refused, so a guard around a deletion cannot pass on an unset
variable.}
\end{center}

\paragraph{The export can still disappear.} After the gate returns, the runner checks that the
export directory is still there. If it is not, the result is reported as its own condition ---
``export lost: another gate's retention removed it while this gate ran; the step failures
above are not a verdict on the tree'' --- rather than as a failure of the code. A gate that
reports the wrong cause is worse than a gate that reports nothing.

\section{What a step is allowed to say}

\texttt{scripts/gate-lib.sh} runs each step in its own session, streams its output, and prints
a heartbeat at least every 30 seconds with the step's last output line. Three verdicts exist:

\begin{itemize}
	\item \texttt{PASS: <label>} --- the command exited 0 and printed no \texttt{SKIPPED:}
	line.
	\item \texttt{SKIPPED: <reason>} --- the command exited 0 and said it skipped. There is no
	PASS line; the skip is the record. \texttt{QG\_SKIP\_COVERAGE=1} uses this.
	\item \texttt{FAIL: <label>} --- non-zero, or killed by a signal, or hung. A step whose
	output has not changed for \texttt{QG\_STALL\_SECS} (default 1\,800) is killed and reported
	as \texttt{HUNG} with the seconds since its last heartbeat and its last line.
\end{itemize}

The heartbeat is a progress note only. The exit code is the sole verdict, which is why the
stall detector exists: a step that prints nothing for thirty minutes is indistinguishable from
a step that is working, and the difference has to be decided by a bound.

\begin{center}
\resizebox{\linewidth}{!}{%
\begin{tikzpicture}[node distance=8mm and 12mm]
	\node[b] (start) {run\_step};
	\node[b, right=of start] (child) {step process\\\scriptsize own session};
	\node[b, below=of child] (hb) {heartbeat\\\scriptsize 30 s};
	\node[v, right=of child] (stall) {output changed?};
	\node[b, below=of stall] (kill) {kill\\\scriptsize HUNG};
	\node[b, above=of stall] (rc) {exit code};
	\node[b, left=of rc, xshift=-2mm] (verdict) {PASS / SKIPPED / FAIL};

	\draw[f] (start) -- (child);
	\draw[f] (child) -- (hb);
	\draw[f] (child) -- (stall);
	\draw[f] (stall) -- node[p, right=2pt, midway]{no, 1800 s} (kill);
	\draw[f] (child) -- (rc);
	\draw[f] (rc) -- (verdict);
\end{tikzpicture}}
\captionof{figure}{Liveness and verdict are separate signals, so a step can be alive and wrong
or dead and therefore not a verdict at all.}
\end{center}

\section{Controls}

The single most useful rule in this repository is that every rule which suppresses, skips or
special-cases something needs a control that must still take the normal path. A suppression
without a control cannot be distinguished from a broken check.

\begin{longtable}{@{}p{0.34\textwidth}p{0.28\textwidth}p{0.30\textwidth}@{}}
\toprule
\textbf{The rule} & \textbf{The near-miss that must still happen} & \textbf{Where it is pinned} \\
\midrule
\endhead
a running gate's export is not reaped & a dead owner's export \emph{is} reaped &
\texttt{fresh\_gate.rs} tests \\
the test step's target lies inside the export & every other step keeps the warm cache &
\texttt{gate\_env} test \\
the sandbox refuses an out-of-grant write & an in-grant write must \emph{succeed} &
\texttt{sandbox.rs} probes \\
a failure report echoes the panic block & a build error with no test failure yields nothing &
\texttt{failure\_report} test \\
a graph refuses an unreachable callee & the reachable callee is still called & splice tests \\
a dirty tree stops dispatch & a row that owns its own dirty tree may resume & lane tests \\
\bottomrule
\end{longtable}

The sandbox pair is the clearest illustration of why the rule exists. A ruleset whose every
rule was skipped passes an out-of-grant probe while granting nothing; every later refusal is
then a coincidence that looks like policy. The in-grant probe is what turns ``something is
denied'' into ``this is granted and that is not''.

\section{A criterion that is already green measures nothing}

An exit criterion is a shell command that fails today. If it passes at \texttt{HEAD}, before
any work, it is not a criterion; it is a wish, and the work it describes may already have
landed or may be described by something else entirely. The rule that follows is mechanical:
every new check is run against \texttt{HEAD} first, and if it passes it is rewritten or
withdrawn.

\begin{center}
\resizebox{\linewidth}{!}{%
\begin{tikzpicture}[node distance=8mm and 12mm]
	\node[b] (red) {run at HEAD\\\scriptsize must FAIL};
	\node[b, right=of red] (impl) {implement};
	\node[b, right=of impl] (green) {run again\\\scriptsize must PASS};
	\node[b, below=of green] (mutate) {mutate the fix back};
	\node[b, left=of mutate] (red2) {run again\\\scriptsize must FAIL};
	\node[b, left=of red2] (restore) {restore};

	\draw[f] (red) -- (impl);
	\draw[f] (impl) -- (green);
	\draw[f] (green) -- (mutate);
	\draw[f] (mutate) -- (red2);
	\draw[f] (red2) -- (restore);
	\draw[bk] (red.south) -- ++(0,-5mm) node[p, below=2pt, midway]{passes: no criterion};
\end{tikzpicture}}
\captionof{figure}{Four runs, not two: the middle pair proves the work did something and the
last pair proves the test is about the work rather than about the tree being green that day.}
\end{center}

The fourth step is the one that gets skipped, and it is the one that catches the most. A test
can pass for a reason unrelated to the fix --- another change in the same commit, a fixture
that happened to be left behind, a guard that fires on its own control. Mutating the fix back
and watching the test go red is the only evidence that the test is about the fix. It is also
how a partly-guarded fix is discovered: mutate one of three guards and see whether anything
notices.

\paragraph{Where this shows up in the store.} The retention work is the clearest example in
this manual. The criterion was a test that seeds an expired run and asserts it is gone, plus
its opposite: a seeded \texttt{Running} run that must survive. At \texttt{HEAD} the first
failed because there was no retention code at all --- the search for \texttt{retention|prune}
in \texttt{pii-store/src} returned nothing. After the implementation both passed. The mutation
that matters is deleting the live-run filter: the first test still passes, the second goes
red. A test suite with only the first would have certified a sweep that can delete work in
flight.

\section{The three lints that encode a convention}

Three gate steps exist because a convention that lives only in a document is a convention that
will be broken. Each encodes a rule the compiler cannot express, and each is a step only
because it can fail.

\begin{itemize}
	\item \textbf{\texttt{lint-files}} enforces 450 \emph{effective} lines per file, where
	effective means non-blank and non-comment. The distinction is not pedantry: a clean
extraction that moves code into a new file without deleting anything leaves both files over
	budget, which is exactly the decision that parked the model-role work. The step also runs
	spec and doc hygiene checks.
	\item \textbf{\texttt{lint-scratch}} fails on a test that uses a crate-local \texttt{\_tmp}
	directory it never creates. The sanctioned helper is
	\texttt{pii\_store::testing::tmp\_scratch}, which creates the crate-local directory
	\emph{and} the temp dir inside it, so a test cannot depend on a directory a previous run
	happened to leave. This class cost four separate fixes, because leftover state made the
	repository gate look green while a clean export failed.
	\item \textbf{\texttt{reachability-check}} proves that a declared capability seam still has
	a real, non-test, non-comment caller. It was added after a resolver existed, was tested,
	and was called by nothing: the test exercised the function directly, so it passed, and the
	feature did not work.
\end{itemize}

The three share a shape: the rule is stated in prose somewhere, the prose is not loaded by the
compiler, and so the rule is restated as a step that exits non-zero. A rule that lives only in
a reference document is a rule that will be broken again --- a sentence this repository has
had to learn twice, once about dispatch scripts and once about commit identity.

\section{The coverage floor}

\texttt{coverage-floor.txt} holds one number, \texttt{74.0}, with the measurement that
justified it recorded in its comments: 77.02\% (21\,741 of 28\,229 lines) on 2026-10-10, with
cargo-llvm-cov 0.9.1 over lib, bins and tests, no doctests. The floor is enforced, not
reported:

\begin{itemize}
	\item A missing \texttt{cargo-llvm-cov} or the llvm-tools component is a FAIL, not a skip.
	A tool that is absent cannot measure, and an unmeasured floor is not a floor.
	\item \texttt{QG\_SKIP\_COVERAGE=1} skips only that step and prints a \texttt{SKIPPED:}
	line naming the variable. It is never implied.
	\item A queue commit refuses a gate that passed without measuring: the cycle reads the log
	for the skip line and parks the row as \texttt{gate-coverage-not-measured}. The distinction
	between ``the floor held'' and ``the floor was not measured'' is the whole point.
	\item The floor is raised when coverage rises and never lowered to make a run green.
\end{itemize}

\section{What the gate cannot see}

A gate's limits belong next to its verdict, or the verdict will be over-read. Four are known
and written down:

\begin{enumerate}
	\item \textbf{The \texttt{hal-gpiod} steps are unexercised.} They compile and lint 117
	lines that have never run on hardware, because this host has no \texttt{/dev/gpiochip*}.
	The gate prints that sentence at the end of every run, so a PASS there cannot read as
	hardware coverage.
	\item \textbf{\texttt{net-check} is a text scan.} It misses a symbol built by a macro or
	\texttt{concat!}, an alias, and any non-\texttt{.rs} source; and its allowlist is per-file,
	so anything inside an allowlisted file is unchecked.
	\item \textbf{\texttt{deps-check} reads metadata.} It cannot see a dependency that arrives
	through a \texttt{\#[path]} include or a \texttt{use} of a crate already in the graph.
	\item \textbf{\texttt{doc-drift-check} reads one file.} It checks that backticked paths in
	\texttt{HANDOFF.md} exist, and nothing about whether the prose is true.
\end{enumerate}

The largest limit is not in that list: a green gate is not a correct change. It says the
checks passed, not that the change did what the spec asked. For a landed fix the check is to
mutate it back and watch the test go red, which is the only evidence that the test is about
the fix.

\section{Gates that fired on the case they allowed}

The recurring defect in this codebase is not a missing guard. It is a guard that fires on the
case it exists to permit, while reporting a reason that is not the real one. Six are recorded
in the manuals' own margin, and each one cost hours:

\begin{itemize}
	\item a \textbf{park} that left its working tree dirty, so the dispatcher refused every
	later row for a night;
	\item a \textbf{dirty-tree guard} that stopped the row whose own work it was (row 368 was
	parked by the guard it was written to permit);
	\item \textbf{prose in a phase field} read as a marker, so a dispatch was skipped;
	\item a \textbf{resolver nothing called}, certified as working by a test that hand-wrote
	the input;
	\item a \textbf{sweep} that deleted rows it could not date;
	\item a \textbf{gate} that inherited the caller's environment and failed on a lock the
	caller held.
\end{itemize}

The defence is the same every time: make the check consume what the live path produces, and
keep a control that must still take the normal path.

\section{One document that had drifted}

\texttt{\_docs/02-guides/quality-gate-and-fresh-tree.md} describes the gate as fourteen steps
and the free-space floor as 2\,048 MB. At this revision the list has twenty steps and the
floor is 12\,288 MB. The document is still the best explanation of \emph{why} each early step
exists, and it is wrong about two numbers a reader would act on.

That is not an argument against documentation; it is an argument for the rule that the script
is the source of truth for the step list and the document explains it. When they disagree, the
script wins --- which is exactly why the gate prints its list and its step count at the start
of every fresh run.
